% ═══════════════════════════════════════════════════════════════════
% §5  열역학 상관식 및 유틸리티 함수
% Source: enex_functions.py
% ═══════════════════════════════════════════════════════════════════
\section{열역학 상관식 및 유틸리티 함수}
\label{sec:enex-functions}

본 절은 상위 모델에서 참조되는 \texttt{enex\_functions.py}의 주요 상관식과 알고리즘을 문서화한다.

\subsection{COP 상관식}

\paragraph{공기열원 히트펌프 --- 냉방 모드.}
외기 온도 $T_0$ [°C]의 함수로 표현된 경험적 회귀식:
\begin{equation}\label{eq:cop-ashp-cool}
  \text{COP}_\text{cool} = a_0 + a_1 T_0 + a_2 T_0^2
\end{equation}
계수는 제조사 카탈로그 데이터에서 적합하였다.

\paragraph{공기열원 히트펌프 --- 난방 모드.}
유사한 다항식 형태:
\begin{equation}\label{eq:cop-ashp-heat}
  \text{COP}_\text{heat} = b_0 + b_1 T_0 + b_2 T_0^2
\end{equation}



\subsection{마찰 계수 (Darcy--Weisbach)}

Colebrook--White 음함수 방정식~\cite{colebrook1939}:
\begin{equation}\label{eq:colebrook}
  \frac{1}{\sqrt{f}} = -2\,\log_{10}\!\left(\frac{\varepsilon/D}{3.7} + \frac{2.51}{\text{Re}\,\sqrt{f}}\right)
\end{equation}
Swamee--Jain 초기 추정값에서 출발하여 반복적으로 풀이한다.

\subsection{자연대류 (수직 평판)}

Churchill \& Chu 상관식~\cite{churchill1975}:
\begin{equation}\label{eq:churchill-chu}
  \overline{\text{Nu}}_L = \left[0.825 + \frac{0.387\,\text{Ra}_L^{1/6}}{\left(1 + (0.492/\text{Pr})^{9/16}\right)^{8/27}}\right]^2
\end{equation}

\subsection{강제대류 (강제 교차유동 원통)}

하천수 등 유동 환경 내부의 파이프 외벽 대류 열전달계수 산출을 위해 원통 표면의 Churchill--Bernstein 모델~\cite{churchill1977}을 참조한다:
\begin{equation}\label{eq:churchill-bernstein}
  \overline{\text{Nu}}_D = 0.3 + \frac{0.62\,\text{Re}_D^{1/2}\,\text{Pr}^{1/3}}{\left[ 1 + (0.4/\text{Pr})^{2/3} \right]^{1/4}} \left[ 1 + \left( \frac{\text{Re}_D}{282000} \right)^{5/8} \right]^{4/5}
\end{equation}
이 상관식은 무한히 넓은 프란틀 수와 광범위한 표면 거칠기 조건에서 널리 활용 가능하여 수열원 침수 코일 파이프의 외벽 대류열전달 평가에 적용된다.

\subsection{유한선원 (FLS) g-함수}

보어홀 열교환기의 지중 열응답~\cite{eskilson1987}:
\begin{equation}\label{eq:fls-g}
  g(t) = \frac{1}{2H}\int_0^H \int_0^H
    \frac{\text{erfc}\!\left(\dfrac{\sqrt{(z-z')^2 + r_b^2}}{2\sqrt{\alpha_s\,t}}\right)}{\sqrt{(z-z')^2 + r_b^2}}\;dz'\,dz
\end{equation}
여기서 $H$는 보어홀 깊이, $r_b$는 보어홀 반경, $\alpha_s$는 지중 열확산율이다.

\subsection{TDMA (Thomas 알고리즘)}

삼중대각 시스템 $a_i T_{i-1} + b_i T_i + c_i T_{i+1} = d_i$를 전진 소거 및 후진 대입으로 풀이한다:
\begin{align}
  c'_i &= \begin{cases} c_0/b_0 & i = 0 \\ c_i\,/\,(b_i - a_i\,c'_{i-1}) & i > 0 \end{cases} \label{eq:tdma-c}\\
  d'_i &= \begin{cases} d_0/b_0 & i = 0 \\ (d_i - a_i\,d'_{i-1})\,/\,(b_i - a_i\,c'_{i-1}) & i > 0 \end{cases} \label{eq:tdma-d}\\
  T_i  &= \begin{cases} d'_{N-1} & i = N-1 \\ d'_i - c'_i\,T_{i+1} & i < N-1 \end{cases} \label{eq:tdma-T}
\end{align}
